\documentclass[10pt,twocolumn]{article}
\usepackage[a4paper,margin=0.72in,columnsep=0.28in]{geometry}
\usepackage{cite,amsmath,amssymb,amsfonts,graphicx,textcomp,xcolor,booktabs,url,times}
\usepackage{titlesec}
\titleformat{\section}{\normalfont\bfseries\scshape}{\thesection.}{0.5em}{}
\titleformat{\subsection}{\normalfont\itshape}{\thesubsection.}{0.5em}{}
\providecommand{\IEEEoverridecommandlockouts}{}
\newenvironment{IEEEkeywords}{\vspace{2pt}\noindent\small\textbf{Index Terms---}}{\par\vspace{4pt}}
\providecommand{\BibTeX}{BibTeX}
\begin{document}
\twocolumn[\begin{center}{\Large\bfseries A Signal-Agnostic Ingestion and Radiomic Feature-Extraction Pipeline for a Multi-Family Clinical Decision-Support System\par}\vspace{6pt}{\normalsize First Author\textsuperscript{1}, Second Author\textsuperscript{2}\par}\vspace{3pt}{\small Graduate Program in Applied Informatics / Biomedical Engineering, University of Fortaleza, Fortaleza, Cear\'a, Brazil\par}{\small Email: ias.pesquisas2026@gmail.com\par}\vspace{12pt}\end{center}]

\begin{abstract}
Clinical decision-support systems (CDSS) are usually built around a single
data modality, which forces institutions to maintain a separate software
stack for every family of biomedical signal. We describe \emph{BioStatusIA},
an open, modular pipeline that ingests heterogeneous biomedical data through a
single entry point and adapts its processing to whatever is technically
feasible. The system automatically classifies an arbitrary file or folder into
one of nine operating modes and routes it to one of three physical signal
families---one-dimensional physiological time series (F1: ECG, EEG, EMG, EOG,
PPG, spirometry), two-dimensional DICOM images (F3: radiography, mammography,
ultrasound), and three-dimensional volumes (F4: CT, MR, PET/SPECT)---plus a
tabular branch for pre-computed clinical features. A single normalized data
contract, \texttt{SinalNormalizado}, decouples the readers from the
domain-specific feature extractors, so that morphological, textural (GLCM),
spectral (Welch PSD) and radiomic biomarkers are computed through a uniform
interface regardless of the input. We report the behaviour of the pipeline on
a real breast-ultrasound dataset (780 images), a clinical tabular dataset, and
synthetic exemplars of each family, and we critically discuss the engineering
trade-offs of a signal-agnostic architecture. The pipeline is released as pure,
side-effect-free functions to encourage reproducibility and reuse.
\end{abstract}

\begin{IEEEkeywords}
clinical decision support, biomedical signal processing, radiomics, feature
extraction, DICOM, medical imaging, software architecture, reproducibility.
\end{IEEEkeywords}

\section{Introduction}
Modern clinical practice generates an unusually broad spectrum of digital
signals: an electrocardiogram is a one-dimensional voltage series, a chest
radiograph is a two-dimensional pixel array stored in DICOM, a computed-tomography
study is a three-dimensional voxel volume, and a laboratory panel is a row of
scalar features. Each of these modalities has, historically, its own analysis
software, its own file formats, and its own community of tools
\cite{gramfort2013,mason2011pydicom,brett2020nibabel}. For a small research
group---such as a graduate laboratory building a proof-of-concept CDSS---this
fragmentation is a serious obstacle: every new modality means a new ingestion
path, a new set of dependencies, and a new opportunity for silent failure.

The central engineering hypothesis of this work is that a large fraction of the
early processing stages---file typing, reading, normalization, statistical
characterization, and biomarker extraction---can be expressed \emph{once} behind
a common data contract, and that the parts that genuinely differ between
modalities can be isolated in thin, replaceable extractors. We call the
resulting system \emph{BioStatusIA}. Its design principle is deliberately
permissive: \emph{the system accepts any input in scope, always runs a
statistical analysis on whatever it receives, and adapts the remainder of the
pipeline to what is feasible}. Supervised model training is attempted only when
labels and a sufficient number of examples are present; otherwise the system
still returns a full descriptive characterization.

This paper concentrates on the \emph{ingestion and feature-extraction}
subsystem---the part of BioStatusIA that turns a raw, possibly unlabelled,
possibly unknown file into a structured vector of biomarkers. The
classification, model-selection and reporting layers are the subject of a
companion paper. Our contributions are: (i) an automatic, filesystem-level
modality detector with nine operating modes; (ii) a single normalized data
contract that decouples readers from extractors across three physical signal
families and a tabular branch; (iii) a set of domain-specific but uniformly
exposed biomarker extractors; and (iv) an honest empirical and architectural
evaluation, including the limitations that a signal-agnostic design imposes.

\section{Related Work}
Radiomics---the high-throughput extraction of quantitative descriptors from
medical images---has become a standard bridge between raw pixels and machine
learning \cite{gillies2016radiomics,haralick1973}. Mature libraries exist for
individual modalities: MNE-Python and WFDB for physiological time series
\cite{gramfort2013}, pydicom for DICOM \cite{mason2011pydicom}, and NiBabel and
SimpleITK for volumetric formats \cite{brett2020nibabel}. General-purpose
scientific stacks such as scikit-image and OpenCV provide the texture and
morphology primitives \cite{vanderwalt2014skimage,pedregosa2011sklearn}.
What is comparatively rare in the literature is an integration layer that treats
these modalities as interchangeable inputs to a \emph{single} decision-support
front end. Commercial PACS and vendor CDSS products are typically
modality-locked, and open academic systems tend to target one signal family at
a time. BioStatusIA occupies this gap: rather than proposing a new descriptor,
it proposes an \emph{architecture} in which existing descriptors are exposed
through one contract and one automatic dispatcher.

\section{System Architecture}
\subsection{Design constraints}
The pipeline is written as a collection of \emph{pure functions}: each stage
receives data and returns data, with no hidden input/output. All persistence
(database writes, HTML rendering, file handling) is concentrated in the web
layer, so the scientific core can be unit-tested and reused in isolation. This
separation is the main reason the same extraction code can be driven both by an
interactive web server and by a batch command-line entry point.

\subsection{Automatic modality detection}
The dispatcher, \texttt{detectar\_estrutura}, receives a path---which may be a
single file, a folder, or an extracted archive---and returns one of nine
in-scope operating modes. For a single file it tests, in order, the predicates
for common images, tabular files, physiological time series, single-slice DICOM,
and 3-D volumes. For a directory it recursively inventories the file types
present and resolves multimodal combinations; folders whose sub-directories
carry recognizable class names (e.g.\ \texttt{benign/}, \texttt{malignant/}) are
promoted to a \emph{labelled dataset} mode. A folder of DICOM files is
disambiguated by cardinality: ten or more slices are read as a 3-D series,
fewer as a set of 2-D images. Inputs outside the supported scope resolve to an
explicit \texttt{invalid} mode rather than failing silently. Table~\ref{tab:modes}
summarizes the modes and the physical family each triggers.

\begin{table}[t]
\centering
\caption{The nine in-scope operating modes and the family they route to.}
\label{tab:modes}
\small
\begin{tabular}{@{}lll@{}}
\toprule
Mode & Trigger & Family \\
\midrule
\texttt{imagem\_unica}      & one raster image            & Image \\
\texttt{imagens\_soltas}    & unlabelled image folder     & Image \\
\texttt{dataset\_rotulado}  & \texttt{benign/}+\texttt{malignant/} & Image \\
\texttt{tabular}            & CSV/TSV/TXT                 & Tabular \\
\texttt{multimodal}         & images + tabular            & Image+Tab \\
\texttt{sinal\_temporal}    & \texttt{.dat/.edf/.mat/...}  & F1 \\
\texttt{imagem\_dicom\_2d}  & single \texttt{.dcm}        & F3 \\
\texttt{volume\_3d}         & \texttt{.nii/.mha}, $\geq$10 \texttt{.dcm} & F4 \\
\texttt{multimodal\_expandido} & mixed families           & F1+F3+F4+Tab \\
\bottomrule
\end{tabular}
\end{table}

\subsection{The normalized data contract}
The keystone of the design is a single dataclass, \texttt{SinalNormalizado},
returned by every reader:
\begin{quote}\small
\texttt{familia}, \texttt{tipo}, \texttt{dados} (\emph{ndarray}),
\texttt{taxa\_amostragem}, \texttt{canais}, \texttt{metadados},
\texttt{caminho\_original}, \texttt{dados\_viz} (a downsampled trace of
$\leq$2000 points for the browser).
\end{quote}
All reading passes through a universal dispatcher, \texttt{carregar\_sinal},
which selects the correct reader by extension---WFDB for
\texttt{.dat/.hea}, MNE for \texttt{.edf/.bdf}, SciPy for \texttt{.mat}, an XML
parser for spirometry exports, pydicom for \texttt{.dcm}, and NiBabel/SimpleITK
for volumes---and returns the same structure in every case. Because the
extractors consume only \texttt{SinalNormalizado}, adding a new file format
requires writing one reader and touching nothing downstream. This contract is
therefore the explicit boundary between the ``read'' and ``understand'' halves
of the system, and we treat it as inviolable.

\section{Feature Extraction}
Each physical family has a dedicated extractor that consumes the normalized
signal and emits a dictionary of biomarkers grouped by domain.

\subsection{F1 --- physiological time series}
Time-domain descriptors (RMS, mean, standard deviation, skewness, kurtosis,
peak-to-peak amplitude, SNR in dB) are complemented by frequency-domain
descriptors obtained from a Welch power-spectral-density estimate: dominant
frequency, spectral centroid, and normalized band powers. Signal-type-specific
analyses are dispatched from the declared type: for ECG the extractor
band-pass-filters the highest-variance channel (0.5--40~Hz, 4th-order
zero-phase Butterworth), detects R-peaks with a physiologically motivated
minimum inter-beat distance, and derives heart rate and the classical
heart-rate-variability indices RMSSD, SDNN and pNN50; for EEG it computes
relative delta/theta/alpha/beta/gamma band powers and the alpha/beta ratio; for
EMG, envelope RMS and median frequency; for spirometry, FVC, FEV1, the
FEV1/FVC ratio and PEF.

\subsection{F3 --- 2-D DICOM}
On single-slice DICOM the extractor computes nine imaging biomarkers---two
morphological (circularity, solidity from the largest external contour), four
textural from the gray-level co-occurrence matrix (contrast, homogeneity,
energy and a histogram entropy), and three intensity descriptors (SNR,
skewness, kurtosis)---augmented by DICOM-specific quantities such as the
fraction of high-density pixels, mean gradient, modality tag and pixel spacing,
with automatic Hounsfield windowing driven by the header's window
centre/width.

\subsection{F4 --- 3-D volumes}
Volumes are re-oriented to an axial-first convention and characterized by global
intensity statistics (mean, dispersion, median, skewness, kurtosis, 5th/95th
percentiles, count of high-intensity voxels), GLCM texture computed
independently on the three orthogonal planes, and 3-D morphology (lesion volume
in mm\textsuperscript{3}, sphericity, bounding-box axes). A helper renders an
arbitrary axial slice to base64 for visualization without exposing a separate
image route.

\subsection{Common image radiomics and the tabular branch}
For ordinary rasters the same nine morphological/textural/intensity biomarkers
as F3 are computed after an adaptive pre-processing stage that chooses denoising,
normalization and equalization from a lightweight analysis of the base
(intensity, contrast, outliers, normality and size consistency). The tabular
branch parses arbitrary CSV/TSV/TXT files: it sniffs the delimiter, tolerates
UTF-8 and Latin-1 encodings, and detects the label column either by name
(\texttt{label}, \texttt{class}, \texttt{diagnosis}, \texttt{target}, \dots) or,
failing that, as the last column with between two and ten distinct values;
binary labels are mapped to $\{0,1\}$ through a keyword table.

\section{Experimental Setup and Results}
We exercised the pipeline on three kinds of input to characterize its behaviour
rather than to claim state-of-the-art accuracy. All experiments use Python~3.11,
scikit-learn, scikit-image, OpenCV, SciPy, MNE, pydicom and NiBabel; random
seeds are fixed for reproducibility.

\subsection{Real breast-ultrasound imaging (BUSI)}
On the Breast Ultrasound Images dataset \cite{aldhabyani2020busi}, the pipeline
autonomously analyzed the base of 780 images (estimated mean intensity 83.6,
mean contrast 51.5, Shapiro $p{=}0.021$, inconsistent sizes) and selected a
Gaussian-denoising / min--max-normalization / 256$\times$256 strategy. Biomarker
extraction succeeded on all 780 images (570 benign, 210 malignant). The
extraction stage is the load-bearing part evaluated here; the downstream
six-model comparison (reported in full in the companion paper) yields, for the
best model, an accuracy of 0.814 and an ROC-AUC of 0.800, with a recall of
0.452 that transparently reflects the untreated class imbalance of this branch.
We regard these realistic---not perfect---numbers as the honest operating point
of the extraction pipeline on genuine clinical data.

\subsection{Clinical tabular data}
On a 50-row breast-cancer panel with 30 numerical features, the tabular parser
correctly identified the \texttt{diagnosis} label, mapped B/M to $0/1$, produced
a balanced 25/25 split and computed univariate statistics for all features
(e.g.\ mean radius $15.32\pm3.90$, mean area $770.6\pm389.8$). The very small
sample size makes any downstream accuracy figure optimistic, a point we
emphasize rather than hide.

\subsection{Synthetic exemplars of each family}
To smoke-test the F1/F3/F4 readers and extractors end-to-end we generated small
synthetic sets (14 exemplars each). The dispatcher correctly identified a
MATLAB ECG signal and produced time and frequency feature groups; a synthetic CT
DICOM produced morphological, GLCM, intensity and DICOM-radiomic groups
(15-dimensional vector); and a synthetic NIfTI volume produced global, 3-D
texture and 3-D morphology groups (13-dimensional vector). These runs confirm
the plumbing of every family but, being synthetic and tiny, are explicitly
\emph{not} evidence of diagnostic performance.

\section{Discussion}
The results support the architectural claim: a single contract and a single
dispatcher are sufficient to carry four data families from raw file to
structured biomarkers, and the pure-function decomposition makes each stage
independently testable. The design also exposes clear limitations. First, a
signal-agnostic front end trades depth for breadth: the GLCM descriptor is
computed at a single distance and orientation and is therefore not
rotation-invariant, and the physiological analyses use robust but basic
estimators rather than modality-optimized toolchains. Second, the extractors
emit vectors of heterogeneous dimensionality (nine to fifteen features across
families), so cross-family model reuse requires an explicit alignment step.
Third, the permissive philosophy that lets the system ``always run something''
can mask degenerate inputs; the single-image smoke test, for instance, produces
biomarkers that are individually valid but statistically meaningless. We view
these not as defects to conceal but as the natural cost of generality, and as a
roadmap: multi-orientation Haralick features, a per-family feature schema, and
input-quality gating are concrete next steps.

\section{Ethical Considerations}
BioStatusIA is a decision-\emph{support} tool. Every visual output and every
generated report carries a mandatory notice that its results do not replace
assessment by a licensed physician. The extraction pipeline computes descriptive
biomarkers only; it makes no autonomous clinical decision, and no result should
be acted upon without professional review. Datasets used here are public and
de-identified.

\section{Conclusion}
We presented the ingestion and feature-extraction core of BioStatusIA, a
signal-agnostic CDSS pipeline that detects its input automatically, normalizes
four data families behind one contract, and extracts domain-appropriate
biomarkers through a uniform interface. Empirical runs on real imaging, real
tabular data and synthetic signals demonstrate that the architecture is sound
and reproducible, while an honest discussion of its limitations frames the path
toward richer descriptors and stronger input validation. The companion paper
builds the clinically-enriched AutoML and multi-agent reporting layers on top of
this foundation.

\begin{thebibliography}{00}
\bibitem{gramfort2013} A. Gramfort \emph{et al.}, ``MEG and EEG data analysis
with MNE-Python,'' \emph{Frontiers in Neuroscience}, vol.~7, art.~267, 2013.
\bibitem{mason2011pydicom} D. Mason, ``SU-E-T-33: Pydicom: an open source DICOM
library,'' \emph{Medical Physics}, vol.~38, no.~6, p.~3493, 2011.
\bibitem{brett2020nibabel} M. Brett \emph{et al.}, ``nibabel: Access a
multitude of neuroimaging data formats,'' Zenodo, 2020.
\bibitem{gillies2016radiomics} R. J. Gillies, P. E. Kinahan, and H. Hricak,
``Radiomics: images are more than pictures, they are data,'' \emph{Radiology},
vol.~278, no.~2, pp.~563--577, 2016.
\bibitem{haralick1973} R. M. Haralick, K. Shanmugam, and I. Dinstein,
``Textural features for image classification,'' \emph{IEEE Trans. Syst., Man,
Cybern.}, vol.~SMC-3, no.~6, pp.~610--621, 1973.
\bibitem{vanderwalt2014skimage} S. van der Walt \emph{et al.},
``scikit-image: image processing in Python,'' \emph{PeerJ}, vol.~2, p.~e453,
2014.
\bibitem{pedregosa2011sklearn} F. Pedregosa \emph{et al.}, ``Scikit-learn:
Machine learning in Python,'' \emph{J. Mach. Learn. Res.}, vol.~12,
pp.~2825--2830, 2011.
\bibitem{aldhabyani2020busi} W. Al-Dhabyani, M. Gomaa, H. Khaled, and A. Fahmy,
``Dataset of breast ultrasound images,'' \emph{Data in Brief}, vol.~28,
p.~104863, 2020.
\end{thebibliography}

\end{document}
